\section*{Теория}
\addcontentsline{toc}{section}{Теория}

\subsection*{0. Базовая теория}

\begin{definition}
    \textbf{Векторное (линейное) пространство} над полем $\mathbb{F}$ (обычно $\mathbb{F} = \mathbb{R}$ или $\mathbb{C}$) — это множество $V$, элементы которого называются \textbf{векторами}, снабжённое двумя операциями: сложением векторов и умножением вектора на скаляр (элемент поля). Эти операции должны удовлетворять следующим аксиомам.

\textit{Аксиомы сложения:}

\begin{enumerate}
    \item \textbf{Коммутативность:} $x + y = y + x$ для любых $x, y \in V$.
    \item \textbf{Ассоциативность:} $(x + y) + z = x + (y + z)$ для любых $x, y, z \in V$.
    \item \textbf{Существование нулевого вектора:} существует элемент $0 \in V$ такой, что $x + 0 = x$ для любого $x \in V$.
    \item \textbf{Существование противоположного вектора:} для каждого $x \in V$ существует элемент $-x \in V$ такой, что $x + (-x) = 0$.
\end{enumerate}

\textit{Аксиомы умножения на скаляр:}

\begin{enumerate}
    \item \textbf{Дистрибутивность умножения на скаляр относительно сложения векторов:} $\lambda(x + y) = \lambda x + \lambda y$ для любых $\lambda \in \mathbb{F}, x, y \in V$.
    \item \textbf{Дистрибутивность относительно сложения скаляров:} $(\lambda + \mu)x = \lambda x + \mu x$ для любых $\lambda, \mu \in \mathbb{F}, x \in V$.
    \item \textbf{Ассоциативность умножения на скаляр:} $(\lambda\mu)x = \lambda(\mu x)$ для любых $\lambda, \mu \in \mathbb{F}, x \in V$.
    \item \textbf{Унитарность:} $1 \cdot x = x$ для любого $x \in V$, где $1$ — единица поля $\mathbb{F}$.
\end{enumerate}
\end{definition}

\begin{definition}
\textbf{Линейный оператор} (или линейное преобразование) пространства $V$ — это отображение $\varphi : V \to V$, удовлетворяющее для любых векторов $x, y \in V$ и любого скаляра $\lambda \in \mathbb{F}$ условиям:
$$ \varphi(x + y) = \varphi(x) + \varphi(y), \quad \varphi(\lambda x) = \lambda\varphi(x). $$
\end{definition}

\begin{definition}
    Множество векторов $e_1, \dots, e_n$ линейного пространства $L$ является его \textit{базисом}, если:
\begin{enumerate}
    \item[1)] Любой вектор $x \in L$ выражается в виде линейной комбинации векторов $e_1, \dots, e_n$;
    \item[2)] Вектора $e_1, \dots, e_n$ линейно независимы.
\end{enumerate}
\end{definition}

Базис векторного пространства $V$ может быть выбран разными способами, но количество векторов в любом базисе $V$ одинаково. Это количество называется \textbf{размерностью} векторного пространства $V$ и обозначается $\dim V$. Если это число конечно, пространство называется \textbf{конечномерным}.

\subsection*{1. Линейные отображения и матрицы}

Пусть $\dim V = n$ и фиксирован базис $\mathcal{E} = (e_1,\dots,e_n)$ в $V$. 

\noindent\textbf{Координатным столбцом} вектора $v = \sum_{i=1}^n v_i e_i$ называется
\[
[v]_\mathcal{E} = \begin{pmatrix}
v_1 \\
v_2 \\
\vdots \\
v_n
\end{pmatrix} \in \mathbb{F}^n.
\]

Отображение $v \mapsto [v]_\mathcal{E}$ является изоморфизмом векторных пространств: оно взаимно однозначно и сохраняет операции. Таким образом, после фиксации базиса абстрактное пространство $V$ становится неотличимым от пространства столбцов $\mathbb{F}^n$.

Теперь рассмотрим линейный оператор $\varphi : V \to V$ и произвольный элемент $v \in V$. Зафиксируем базис $\mathcal{E}$ в $V$. Тогда
\[
[\varphi(v)]_\mathcal{E} = \left[\varphi\left(\sum_{i=1}^n v_i e_i\right)\right]_\mathcal{E} = \sum_{i=1}^n [\varphi(e_i)]_\mathcal{E} v_i = ([\varphi(e_1)]_\mathcal{E}\ [\varphi(e_2)]_\mathcal{E}\ \dots\ [\varphi(e_n)]_\mathcal{E}) \cdot [v]_\mathcal{E}.
\]

\textbf{Матрицей оператора} $\varphi$ в базисе $\mathcal{E}$ называется квадратная матрица $M_\mathcal{E}(\varphi)$ размера $n \times n$, $j$-й столбец которой есть координатный столбец образа $j$-го базисного вектора $e_j$ в том же базисе $\mathcal{E}$:
\[
M_\mathcal{E}(\varphi) = ([\varphi(e_1)]_\mathcal{E}\ [\varphi(e_2)]_\mathcal{E}\ \dots\ [\varphi(e_n)]_\mathcal{E}).
\]

Как показано выше, для любого вектора $v \in V$ выполнено соотношение
\[
[\varphi(v)]_\mathcal{E} = M_\mathcal{E}(\varphi) \cdot [v]_\mathcal{E}.
\]

То есть, после отождествления векторов с их координатными столбцами, действие абстрактного линейного оператора $\varphi$ превращается в умножение слева на конкретную числовую матрицу. Обратно, любая квадратная матрица $A$ задаёт линейный оператор $[v]_\mathcal{E} \mapsto A[v]_\mathcal{E}$. Это отождествление является краеугольным камнем линейной алгебры.

\subsection*{2. Матрица перехода и смена базиса}

В разных базисах одному и тому же линейному оператору могут соответствовать разные матрицы. Связаны эти матрицы будут с помощью так называемой матрицы перехода, которая зависит от этих базисов, но не от оператора. То есть при смене базиса матрицы всех линейных операторов меняются «одинаковым образом».

Пусть в пространстве $V$ заданы два базиса: старый $\mathcal{E} = (e_1, \dots, e_n)$ и новый $\mathcal{F} = (f_1, \dots, f_n)$. \textbf{Матрицей перехода} от $\mathcal{E}$ к $\mathcal{F}$ называется матрица $P_{\mathcal{E} \to \mathcal{F}}$, столбцы которой суть координаты новых базисных векторов $f_j$ в старом базисе $\mathcal{E}$:
\[
(f_1, \dots, f_n) = (e_1, \dots, e_n)P_{\mathcal{E} \to \mathcal{F}},
\]
\[
P_{\mathcal{E} \to \mathcal{F}} = ([f_1]_\mathcal{E}\ [f_2]_\mathcal{E}\ \dots\ [f_n]_\mathcal{E}).
\]
Эта матрица невырождена (обратима), поскольку базисы линейно независимы.

\textit{Преобразование координат:} Если $[v]_\mathcal{E}$ — координаты вектора $v$ в старом базисе, а $[v]_\mathcal{F}$ — в новом, то
\[
[v]_\mathcal{E} = P_{\mathcal{E} \to \mathcal{F}} \cdot [v]_\mathcal{F}, \quad [v]_\mathcal{F} = P_{\mathcal{E} \to \mathcal{F}}^{-1} \cdot [v]_\mathcal{E}.
\]

\vspace{1em}
\textit{Смена матрицы линейного оператора:} Пусть $\varphi : V \to V$ — линейный оператор. В старом базисе $\mathcal{E}$ его матрица равна $M_\mathcal{E}(\varphi)$, в новом базисе $\mathcal{F}$ — $M_\mathcal{F}(\varphi)$. Тогда выполняется соотношение
\[
M_\mathcal{F}(\varphi) = P_{\mathcal{E} \to \mathcal{F}}^{-1} \cdot M_\mathcal{E}(\varphi) \cdot P_{\mathcal{E} \to \mathcal{F}}.
\]

Матрицы $A, B \in M_n(\mathbb{F})$, связанные соотношением $B = C^{-1}AC$ для некоторой невырожденной матрицы $C \in M_n(\mathbb{F})$ (то есть отличающиеся заменой базиса), называются \textbf{подобными}.

\subsection*{3. Спектральная теория и диагонализация}

\begin{definition}
Пусть $L$ — линейное пространство над полем $\mathbb{F}$ (чаще всего $\mathbb{R}$ или $\mathbb{C}$), $\varphi : L \to L$ — линейный оператор. Ненулевой вектор $x$ называется \textbf{собственным вектором} оператора $\varphi$, если $\varphi(x) = \lambda x$ для некоторого $\lambda \in \mathbb{F}$. При этом число $\lambda$ называется \textbf{собственным значением} оператора $\varphi$.
\end{definition}

\begin{definition}
Множество собственных векторов, соответствующих одному собственному значению $\lambda$ оператора $\varphi$, дополненное нулевым вектором, образует подпространство, которое называют \textbf{собственным подпространством} оператора $\varphi$ для собственного значения $\lambda$ и обозначают $E_\lambda$:
\[
E_\lambda = \{x \in L \mid \varphi(x) = \lambda x\} = \ker(\varphi - \lambda \mathrm{id}).
\]
Размерность $\dim E_\lambda$ называется \textbf{геометрической кратностью} собственного значения $\lambda$. В матричной форме она вычисляется как $n - \mathrm{rank}(A - \lambda E)$, где $A$ — матрица оператора.
\end{definition}

\begin{definition}
\textbf{Характеристическим многочленом} квадратной матрицы $A$ называется
\[
\chi_A(\lambda) = \det(A - \lambda E).
\]
Далее мы отождествляем понятия линейного оператора в конечномерном пространстве и его матрицы, подразумевая, что мы фиксировали некий базис и задали оператор матрицей. Число $\lambda_0$ является собственным значением оператора тогда и только тогда, когда $\lambda_0$ является корнем характеристического многочлена его матрицы. Кратность $\lambda_0$ как корня многочлена $\chi_A(\lambda)$ называется его \textbf{алгебраической кратностью}.
\end{definition}

\noindent\textbf{Замечание.} Пусть $A$ — квадратная матрица порядка $n$, $\lambda$ — собственное значение $A$. Тогда геометрическая кратность $\lambda$ есть $\dim E_\lambda = n - \mathrm{rk}(A - \lambda E)$.

\noindent\textbf{Утверждение.} Геометрическая кратность собственного значения $\lambda$ не превосходит алгебраической кратности собственного значения $\lambda$.

\noindent\textbf{Теорема.} Следующие четыре условия эквивалентны

\begin{enumerate}
    \item[1)] Матрица $A$ диагонализируема (то есть существует базис, в котором матрица $A'$ оператора, который в данном базисе задан матрицей $A$, диагональна).
    \item[2)] Существует базис линейного пространства, состоящий из собственных векторов матрицы $A$.
    \item[3)] Сумма геометрических кратностей всех собственных значений матрицы $A$ равна размерности линейного пространства (в частности, это заведомо выполнено, когда матрица имеет $n$ различных собственных значений, где $n$ — размерность линейного пространства).
    \item[4)] $\chi_A(\lambda)$ раскладывается на линейные множители над $\mathbb{F}$ и геометрические кратности собственных значений совпадают с алгебраическими кратностями.
\end{enumerate}

\noindent\textbf{Следствие.} Если матрица $A$ порядка $n$ имеет $n$ попарно различных собственных значений, то она диагонализируема.

\noindent\textbf{Следствие.} Если матрица $A$ диагонализируема, то существует невырожденная матрица $C$, столбцы которой состоят из собственных векторов матрицы $A$, такая, что
\[
A=CDC^{-1},
\]
где
\[
D=\operatorname{diag}(\lambda_1,\lambda_2,\ldots,\lambda_n)
\]
— диагональная матрица, на главной диагонали которой расположены собственные значения матрицы $A$.

Для любого $k\in\mathbb N$
\[
A^k=CD^kC^{-1},
\]
где
\[
D^k=\operatorname{diag}(\lambda_1^k,\lambda_2^k,\ldots,\lambda_n^k).
\]

\subsection*{4. Ядро и образ линейного отображения}

\begin{definition}
Пусть $\phi : V \to W$ — линейное отображение.

\textbf{Ядром} отображения $\phi$ называется множество
\[
\ker\phi=\{x\in V\mid \phi(x)=0\}.
\]
\end{definition}

Ядро линейного отображения является линейным подпространством пространства $V$.

\begin{definition}
\textbf{Образом} линейного отображения $\phi$ называется множество
\[
\operatorname{Im}\phi=\{\phi(x)\mid x\in V\}.
\]
\end{definition}

Образ линейного отображения является линейным подпространством пространства $W$.

\vspace{1em}

Если отображение $\phi$ задано матрицей $A$, то

\begin{itemize}
    \item ядро состоит из всех решений однородной системы
    \[
    Ax=0;
    \]
    \item образ совпадает с линейной оболочкой столбцов матрицы $A$.
\end{itemize}

\begin{theorem}
Пусть $\phi:V\to W$ — линейное отображение, где $V$ конечномерно. Тогда
\[
\dim\ker\phi+\dim\operatorname{Im}\phi=\dim V.
\]
\end{theorem}

\noindent\textbf{Следствие.}
Линейное отображение
\[
\phi:V\to W
\]
является сюръективным тогда и только тогда, когда
\[
\operatorname{Im}\phi=W,
\]
то есть
\[
\operatorname{rank}A=\dim W.
\]

\subsection*{5. Инвариантные подпространства}

\begin{definition}
Подпространство $U\subseteq V$ называется \textbf{инвариантным} относительно линейного оператора $\phi$, если
\[
\phi(U)\subseteq U.
\]
\end{definition}

В частности, прямая
\[
L=\langle v\rangle
\]
является инвариантной тогда и только тогда, когда
\[
\phi(v)=\lambda v
\]
для некоторого числа $\lambda$, то есть вектор $v$ является собственным вектором оператора.

\subsection*{6. Сопряжённый оператор}

\begin{definition}
Пусть $V$ — евклидово пространство со скалярным произведением $(\cdot,\cdot)$.

Линейный оператор $f^*:V\to V$ называется \textbf{сопряжённым} к оператору $f$, если для любых $x,y\in V$
\[
(f(x),y)=(x,f^*(y)).
\]
\end{definition}

Пусть базис $e_1,\dots,e_n$ пространства $V$ не обязательно ортонормирован. Его \textbf{матрицей Грама} называется матрица
\[
G=((e_i,e_j))_{i,j=1}^n.
\]

Если матрица оператора $f$ в данном базисе равна $A$, то матрица сопряжённого оператора вычисляется по формуле
\[
A_{f^*}=G^{-1}A^TG.
\]

В частности, если базис ортонормирован, то
\[
G=E,
\]
поэтому
\[
A_{f^*}=A^T.
\]
